%%%%%%%%%%%%%%%%%%%%%%%%%%%%%%%%%%%%%%%%%%%%%%%%%%%%%%%%%%%%%%%%%%
\section{Mathematical model}
\label{sec:mathmodel}
%%%%%%%%%%%%%%%%%%%%%%%%%%%%%%%%%%%%%%%%%%%%%%%%%%%%%%%%%%%%%%%%%%

\noindent
The coupled problem consists of a slender flexible structure, represented by a geometrically exact beam, embedded inside an incompressible flow domain. There is no body-fitted mesh around the structure. The beam exchanges momentum with the flow through a distributed volumetric source, and the surrounding Eulerian fluid mesh remains stationary throughout. Boldface symbols denote vector or, where stated, second-order tensor quantities.

%%%%%%%%%%%%%%%%%%%%%%%%%%%%%%%%%%%%%%%%%%%%%%%%%%%%%%%%%%%%%%%%%%
\subsection{Fluid subsystem}
\label{sec:fluidSubsystem}
%%%%%%%%%%%%%%%%%%%%%%%%%%%%%%%%%%%%%%%%%%%%%%%%%%%%%%%%%%%%%%%%%%

\noindent
The fluid is modelled as an incompressible Newtonian flow governed by
\begin{equation}
\nabla \cdot \mathbf{U} = 0,
\label{eq:continuity}
\end{equation}
\begin{equation}
\frac{\partial (\rho \mathbf{U})}{\partial t}
+ \nabla \cdot (\rho \mathbf{U}\mathbf{U})
= -\nabla p
+ \nabla \cdot \left[ \mu \left(\nabla \mathbf{U} + \nabla \mathbf{U}^{T}\right)\right]
+ \rho \mathbf{g}
+ \mathbf{S}_{m},
\label{eq:momentum}
\end{equation}
where $\mathbf{U}$ is the velocity, $\rho$ the density, $\mu$ the dynamic viscosity, $p$ the pressure, and $\mathbf{g}$ the gravitational acceleration. The term $\mathbf{S}_{m}$ is the volumetric momentum source through which the beam acts on the flow, derived in Section~\ref{sec:almProjection}. The coupled pressure--velocity system is solved with the standard segregated PIMPLE procedure of OpenFOAM, and $\mathbf{S}_{m}$ is updated once per time step rather than at every outer iteration of that procedure, as set out in Section~\ref{sec:almParallelSearch}.

\noindent
The coupling operator is formulated independently of the number of fluid phases. The channel-flow benchmarks of Sections~\ref{sec:beamtunnel} and~\ref{sec:leaflet} are single-phase and the wave-flume validation of Section~\ref{sec:flexveg} is two-phase, with the same implementation used unchanged in both. For the two-phase case a volume-of-fluid formulation transports a phase fraction $\alpha$ to distinguish the two immiscible fluids, the system is treated as a single mixture with phase-averaged properties $\rho = \alpha \rho_1 + (1-\alpha)\rho_2$ and $\mu = \alpha \mu_1 + (1-\alpha)\mu_2$, and the phase fraction is advanced by
\begin{equation}
\frac{\partial \alpha}{\partial t}
+ \nabla \cdot (\mathbf{U}\alpha)
+ \nabla \cdot \left(\mathbf{U}_r\,\alpha(1-\alpha)\right) = 0,
\label{eq:alphaTransportCompressed}
\end{equation}
in which the third term is the interface-compression flux that counteracts numerical smearing, with $\mathbf{U}_r$ a compression velocity active only near the interface. Boundedness of $\alpha$ is enforced by the MULES limiter \citep{damian_extended_2014}, and the momentum equation is solved in the dynamic-pressure form $p_{\mathrm{rgh}} = p - \rho\,\mathbf{g}\cdot\mathbf{x}$ standard in the \texttt{interFoam} family of solvers.

\subsubsection{Regular wave theories, generation, and absorption}
\label{sec:theory_waves}
The free-surface case of Section~\ref{sec:flexveg} is driven by regular waves prescribed at an inlet boundary. Progressive waves of permanent form on a finite depth $d$ are described by the Stokes expansion, a perturbation series in the steepness $ka$ with $k = 2\pi/\lambda$ the wavenumber and $a = H/2$ the amplitude of a wave of height $H$, whose frequency $\omega = 2\pi/T$ follows from $\omega^2 = g k \tanh kd$ \citep{dean_water_1991}. The adequacy of a given order is judged by the steepness together with the Ursell number $\mathrm{Ur} = H\lambda^2/d^3$, second-order theory being adequate for $\mathrm{Ur} \lesssim 10$ \citep{dean_water_1991}, so where the simulated conditions approach $\mathrm{Ur} \approx 20$ the fifth-order theory of \citet{fenton_fifth_1985} is used instead. The second harmonic of the expansion sharpens the crests, flattens the troughs, and makes the crest-phase velocity larger in magnitude but shorter in duration than the trough-phase velocity. This asymmetry matters for drag-driven structures, because a quadratic drag force rectifies it into a net forcing asymmetry over the cycle.

\noindent
The chosen theory is imposed at the inlet as time-dependent boundary values of velocity and phase fraction, ramped smoothly from rest, combined with active absorption that cancels waves reflected back toward the inlet \citep{higuera_realistic_2013}. At the outlet, a shallow-water absorption condition removes the outgoing wave using the linear long-wave celerity $\sqrt{gd}$. Because that celerity is exact only in the long-wave limit, the outlet absorption degrades as the relative depth $kd$ approaches unity, which is relevant to one of the validation conditions of Section~\ref{sec:flexveg}.

%%%%%%%%%%%%%%%%%%%%%%%%%%%%%%%%%%%%%%%%%%%%%%%%%%%%%%%%%%%%%%%%%%
\subsection{Beam subsystem}
\label{sec:beamSection}
%%%%%%%%%%%%%%%%%%%%%%%%%%%%%%%%%%%%%%%%%%%%%%%%%%%%%%%%%%%%%%%%%%

\noindent
The structure is modelled as a one-dimensional \emph{Simo--Reissner} beam whose centreline can undergo arbitrarily large three-dimensional translations and rotations without any small-rotation linearisation, and which retains shear deformation and torsion in addition to bending and axial stretch. A \emph{total Lagrangian} description is adopted, in which every kinematic quantity is referred back to the fixed reference configuration through the arc-length coordinate $s \in [0,L]$. The initial stress-free centreline is $\mathbf{r}_0(s)$ and the current deformed centreline is $\mathbf{r}(s,t) = \mathbf{r}_0(s) + \mathbf{w}(s,t)$, and the orientation of the cross-section at each $s$ is represented by a rotation tensor $\boldsymbol{\Lambda}(s,t) \in SO(3)$. Finite rotations are parameterised through a rotation vector and the exponential map $\boldsymbol{\Lambda}(\boldsymbol{\psi}) = \exp(\widehat{\boldsymbol{\psi}})$, with $\widehat{(\cdot)}$ the skew-symmetric matrix representing the cross product, and the Newton iteration updates them multiplicatively so that every iterate stays exactly on $SO(3)$ \citep{dai_eulerrodrigues_2015,spurrier_comment_1978}.

\noindent
The strain measures are the co-rotated translational strain $\boldsymbol{\Gamma} = \boldsymbol{\Lambda}^T \mathbf{r}' - \boldsymbol{\Lambda}_0^T \mathbf{r}_0'$, whose components are axial extension and shear, and the curvature $\mathbf{K}$, whose components are torsion and bending, with $(\cdot)' = \partial(\cdot)/\partial s$ and $\boldsymbol{\Lambda}_0(s)$ the rotation tensor of the initial configuration. Both are objective, since rotating each quantity back into the local body-attached basis removes any superposed rigid-body motion. A hyperelastic model quadratic in the strain measures gives the material resultants $\mathbf{N} = \mathbf{C}_N \boldsymbol{\Gamma}$ and $\mathbf{M} = \mathbf{C}_M \mathbf{K}$, with the diagonal stiffness matrices $\mathbf{C}_N = \mathrm{diag}(EA,\,GA_2,\,GA_3)$ and $\mathbf{C}_M = \mathrm{diag}(GJ,\,EI_2,\,EI_3)$ appropriate for the bi-symmetric cross-sections used in every case reported here, and the spatial resultants follow by push-forward, $\mathbf{n} = \boldsymbol{\Lambda}\mathbf{N}$ and $\mathbf{m} = \boldsymbol{\Lambda}\mathbf{M}$. The kinematics, the strain measures, and the linearisation of the resultants are set out in full by \citet{bali_cell-centered_2022,bali_finite_2024}, and only what the coupling needs is repeated here.

\noindent
The balance of linear and angular momentum along the centreline reads
\begin{equation}
\mathbf{n}' + \mathbf{f} \;=\; \rho_{b}A\,\ddot{\mathbf{w}},
\label{eq:ForceBalanceDynamic}
\end{equation}
\begin{equation}
\mathbf{m}' + \mathbf{r}'\times\mathbf{n} + \mathbf{t} \;=\;
\boldsymbol{\Lambda}\,\bigl(\mathbf{I}_{\rho}\,\dot{\boldsymbol{\Omega}}
\;+\; \boldsymbol{\Omega}\times\bigl(\mathbf{I}_{\rho}\,\boldsymbol{\Omega}\bigr)\bigr),
\label{eq:MomentBalanceDynamic}
\end{equation}
where $\mathbf{f}$ and $\mathbf{t}$ are the external distributed force and torque per unit length, $\rho_{b}A$ is the mass per unit reference length, $\ddot{\mathbf{w}}$ the translational acceleration of the centreline, $\boldsymbol{\Omega}$ the material angular velocity defined through $\widehat{\boldsymbol{\Omega}} = \boldsymbol{\Lambda}^{T}\dot{\boldsymbol{\Lambda}}$, and $\mathbf{I}_{\rho} = \rho_{b}\,\mathrm{diag}\left(J,\;I_{2},\;I_{3}\right)$ the material rotary-inertia tensor. The right-hand side of Eq.~\eqref{eq:MomentBalanceDynamic} is the push-forward of the material angular-momentum rate and contains both the rotary-inertia and the gyroscopic term. Setting the right-hand sides to zero recovers the quasi-static form used for the static benchmark of Section~\ref{sec:cantilever}. No physical structural damping is added in any simulation reported here, and the hydrodynamic drag loading provides the physical damping of the coupled response.

%%%%%%%%%%%%%%%%%%%%%%%%%%%%%%%%%%%%%%%%%%%%%%%%%%%%%%%%%%%%%%%%%%
\subsection{Hydrodynamic closure on the beam}
\label{sec:beamExternalLoading}
%%%%%%%%%%%%%%%%%%%%%%%%%%%%%%%%%%%%%%%%%%%%%%%%%%%%%%%%%%%%%%%%%%

\noindent
The distributed external force per unit length in Eq.~\eqref{eq:ForceBalanceDynamic} comprises a hydrodynamic drag force and a net buoyancy force, $\mathbf{f} = \mathbf{f}_{d} + \mathbf{f}_{b}$. The drag follows the quasi-steady Morison formulation \citep{morison_force_1950}, split into tangential and normal components relative to the local beam tangent,
\begin{equation}
\mathbf{f}_{d} \;=\; \tfrac{1}{2}\,\rho_{f}\,d\,
\bigl(C_{dt}\,\lVert\mathbf{V}_{r,t}\rVert\,\mathbf{V}_{r,t}
\;+\; C_{dn}\,\lVert\mathbf{V}_{r,n}\rVert\,\mathbf{V}_{r,n}\bigr),
\label{eq:fext_drag}
\end{equation}
where $\rho_{f}$ is the fluid density, $d$ the cross-section drag reference length, $C_{dt}$ and $C_{dn}$ the tangential and normal drag coefficients, and $\mathbf{V}_{r,t}$ and $\mathbf{V}_{r,n}$ the tangential and normal projections of the relative velocity between the beam and the surrounding fluid. That surrounding velocity is the upstream-sampled velocity provided by the operator of Section~\ref{sec:almCouplingTheory}, and the discrete form of the line force is given in Section~\ref{sec:almLineForce}. The net effect of gravity and buoyancy on a submerged segment is
\begin{equation}
\mathbf{f}_{b} \;=\; \rho_{b}\,A\,\!\left(\frac{\rho_{b}-\rho_{f}}{\rho_{b}}\right)\mathbf{g}.
\label{eq:fext_buoyancy}
\end{equation}

\noindent
The closure carries no added-mass term. When a beam accelerates through the surrounding fluid, the displaced fluid exerts an additional inertial reaction, which the classical Morison formulation represents as
\begin{equation}
\mathbf{f}_{a} \;=\; \rho_{f}\,A\,\bigl(
C_{Mt}\,\mathbf{a}_{r,t}
\;+\; C_{Mn}\,\mathbf{a}_{r,n}
\;+\; \mathbf{a}_{f}\bigr),
\label{eq:fext_addedmass}
\end{equation}
with $\mathbf{a}_{r} = \mathbf{a}_{b} - \mathbf{a}_{f}$ the relative acceleration, $C_{Mt}$ and $C_{Mn}$ the added-mass coefficients, and the third term the Froude--Krylov contribution carried by the undisturbed fluid acceleration. Equation~\eqref{eq:fext_addedmass} is not evaluated in the present model, so $C_{Mt} = C_{Mn} = 0$ throughout and the closure is quasi-steady and drag-only.

\noindent
Equations~\eqref{eq:fext_drag} and~\eqref{eq:fext_buoyancy} are therefore a modelling approximation of the hydrodynamic load rather than an exact expression for it, since a body-fitted calculation obtains the force by integrating the resolved pressure and viscous traction over the wetted surface and so carries the added-mass and history contributions that this closure omits. That omission is a deliberate scope decision, and its size, together with the range of conditions in which the approximation remains acceptable, is quantified against benchmark and experimental data in Sections~\ref{sec:leaflet} and~\ref{sec:flexveg} rather than assumed.
